\section{Introduction}

The nocturnal Stable Boundary Layer (SBL) remains one of the most persistent challenges in geophysical fluid dynamics due to the highly non-linear, non-equilibrium coupling between turbulent dissipation, surface energy budgets, and mesoscale forcing \citep{nieuwstadt1984, mahrt2014, vanDeWiel2017}. A central unresolved problem in boundary-layer meteorology is the so-called ``Richardson Number Universality Crisis'': the observation that the critical Bulk Richardson number ($Ri_c$), traditionally expected to be a quasi-universal constant near $0.25$, varies wildly across different environmental regimes \citep{Grachev2013, Monahan2015}. While mid-latitude field campaigns such as CASES-99 report collapse thresholds near the classical limit ($Ri \approx 0.2\text{--}0.4$), polar campaigns like SHEBA record active turbulence persisting at stability ratios exceeding unity ($Ri > 1.2$) \citep{poulos2002cases99, Grachev2005}.

\subsection{The Paradox of Stability Thresholds}

For decades, the search for a ``better'' stability function has assumed that the variability in $Ri_c$ stems from inconsistent physics or measurement noise. However, traditional parameterizations---such as Monin--Obukhov Similarity Theory (MOST) or Mellor--Yamada $K$-theory---rely on an instantaneous local equilibrium assumption, $e = e(S, N^2)$, that erases the system's underlying temporal history and structural hysteresis \citep{MellorYamada1982}. This diagnostic approach forces a smooth, single-valued response that cannot account for the abrupt, ``brittle'' transitions between weakly stable and very stable regimes.

\subsection{Thesis: Richardson Thresholds as Projections, Not Invariants}

In this paper, we propose that the Richardson threshold paradox is not evidence of incompatible physics, but rather a deterministic consequence of dimensionality reduction. We frame the SBL as a singularly perturbed dynamical system governed by Geometric Singular Perturbation Theory (GSPT) \citep{Fenichel1979, Kuehn2015}. Within this framework, our central thesis is that \emph{Richardson thresholds are low-dimensional coordinate projections of a higher-dimensional folded equilibrium manifold}.

By exploiting the natural separation of timescales---where fast turbulence adjustment ($\tau_e \sim 10^1\text{--}10^2\,\text{s}$) responds to slow surface thermal and geostrophic evolution ($\tau_{\text{slow}} \sim 10^3\text{--}10^4\,\text{s}$)---we demonstrate that the ``critical'' threshold is not a static point, but a state-dependent fold locus ($\mathcal{C}_{\text{fold}}$) that deforms dynamically in response to surface energy balance \citep{VanDeWiel2012}.

\subsection{The ``Fold Illusion'' and Emergent Coupling}

A key conceptual contribution of this work is the distinction between isolated boundary-driven decay and true fold catastrophes. We prove that while an isolated atmospheric subsystem exhibits only a monotonic boundary-crossing at the laminar floor ($e = 0$), a genuine S-shaped fold only emerges through the non-linear coupling of the atmosphere to the Surface Energy Balance (SEB). We term the historical conflation of these two mechanisms the ``Fold Illusion,'' and show that the geometric resolution of this illusion is necessary to model intermittent bursting and nocturnal relaxation oscillations accurately.

\subsection{Objectives and Scope of Part 1}

As the first in a three-part series, this manuscript establishes the mathematical foundations and observational resolution of the GSPT-SBL framework. We define a hierarchy of models---from 2D analytical baselines to 5D thermodynamic systems---and provide a sequence of structural theorems:
\begin{itemize}
    \item \textbf{Theorem 1 (Existence and Smoothness):} Establishing the critical manifold ($\mathcal{S}_0$) as a smooth embedded submanifold of the state space.
    \item \textbf{Theorem 2 (Fold Characterization):} Defining the geometric boundary where normal hyperbolicity is lost, triggering turbulence collapse.
    \item \textbf{Theorem 3 (Projection Theorem):} Formalizing the mapping from invariant state-space geometry to campaign-specific scalar diagnostics ($Ri_{\text{fold}}$).
\end{itemize}

Part 1 provides the ``geometric skeleton'' required for the data-driven discovery and prognostic parameterizations developed in Parts 2 and 3.